\worksheettitle{M06\(\star\). Следующее по величине}{\modulemark{M06\(\star\)} Усложнение}
\studentnameblock

\begin{notebox}[Идея]
  Чтобы получить второе по величине число, нужно оставить старшие цифры
  максимально большими и сделать самое маленькое возможное уменьшение как
  можно правее.
\end{notebox}

\begin{practicebox}[Разберите пример]
  Из цифр 2, 5, 7, 9 максимум равен $9752$.
  \begin{enumerate}[label=\textbf{\arabic*.}]
    \item Можно ли сохранить первые две цифры 97? \answerblank[20mm].
    \item Поменяйте местами только последние две цифры. Получится
      \answerblank[35mm].
    \item Почему это число второе по величине? \answerblank[72mm].
  \end{enumerate}
\end{practicebox}

\begin{practicebox}[Найдите два лидера]
  \begin{center}
  \renewcommand{\arraystretch}{1.6}
  \begin{tabular}{c|c|c}
    \toprule
    цифры & наибольшее & второе по величине \\
    \midrule
    1, 3, 6 & \answerblank[28mm] & \answerblank[34mm] \\
    0, 2, 5, 8 & \answerblank[28mm] & \answerblank[34mm] \\
    1, 4, 4, 9 & \answerblank[28mm] & \answerblank[34mm] \\
    0, 0, 3, 7 & \answerblank[28mm] & \answerblank[34mm] \\
    \bottomrule
  \end{tabular}
  \end{center}
\end{practicebox}

\begin{warningbox}[Ловушка с одинаковыми цифрами]
  Если поменять местами две одинаковые цифры, число не изменится. Для цифр
  1, 4, 4, 9 найдите перестановку, которая действительно уменьшает максимум
  наименьшим возможным образом: \answerblank[45mm].
\end{warningbox}

\begin{checkpointbox}[Обоснование]
  Для цифр 0, 2, 5, 8 объясните, почему между найденными первым и вторым
  числами нет другого допустимого числа: \answerblank[110mm].
\end{checkpointbox}
